\CatchFileBetweenTags{\BoundaryStorageVal}{constants.tex}{BoundaryStorageVal}
\CatchFileBetweenTags{\GaugeLoadVal}{constants.tex}{GaugeLoadVal}
\CatchFileBetweenTags{\TotalBandwidthVal}{constants.tex}{TotalBandwidthVal}

\section{Synthesis: Conditional Solvency and Joint Closure}
\label{sec:Synthesis}

The general rules constrain resource use, reversible encoding and interface transfer. Applying them to the vacuum additionally requires a substrate-derived workload, response law and selection criterion. This section distinguishes exact accounting identities from the joint tests that can fail.

\subsection{Capacity Accounting}
The proposed three-component ledger defines
\begin{equation}
 B_{res}=\nu-\frac{\chi}{\sigma-\chi}-\alpha,
 \qquad
 B_{res}+\frac{\chi}{\sigma-\chi}+\alpha=\nu.
\end{equation}
The displayed numerical sum is
\begin{equation}
 \ResidualCapVal+\BoundaryStorageVal+\GaugeLoadVal=\TotalBandwidthVal.
\end{equation}
Lean proves the identity and positivity of the stated remainder under the inherited coupling bounds. The equality is exact by definition. It does not independently select the allocations, establish the gravitational attenuation law, or prove that every allocated resource is fully utilized. Any reserve made available to another interface must be accounted for in the same ledger rather than spent twice.

\subsection{Reversible Encoding and Spectral Interpretation}
The general encoding theorem requires metadata to retain distinctions lost by a coarse description and bounds the required metadata alphabet. The identification of the Margin coefficient with an electron Yukawa or a lowest stable mass is an additional spectral and dynamical hypothesis. A finite encoding cost is not itself an infinite heat penalty, and no mass spectrum follows from metadata cardinality alone.

The electron Yukawa and VEV also provide a joint test. The earlier independent formula $PM(1+\tfrac{\sigma}{D}\alpha)$ gives
\begin{equation}
 0.99972687<\frac{y_e v}{\sqrt{2}m_e}<0.99972689,
\end{equation}
independent of the mass unit. The two formulas therefore cannot both be exact. A field rescaling leaves $y_ev/\sqrt2$ unchanged, and holding the Fermi formula fixed pins that rescaling to $\pm1$, so no normalization choice removes the residual (\texttt{RepairSpace}). Since the VEV reproduces the Fermi-constant value of $v$ within its uncertainty, the Yukawa is now taken from the mass law, $y_e=\sqrt2m_e/v$, which closes exactly. An independent Margin dressing would need the coefficient $1.28777<c<1.28779$ (\cref{subsec:margin_electron}).

\subsection{Surface and Core Interface Tests}
The proposed surface comparison gives close, unequal impedances. For its stated positive-real transfer law, the reflected fraction lies between $10^{-9}$ and $10^{-8}$ (\texttt{SurfaceBudget}). This is compatible with sufficiently small reflected-load cost only if a corresponding reserve is actually available. The scalar's maximal-gain ledger leaves no reserve of its own, so maximal gain cannot pay for its reflection; exact matching leaves unused slack. The largest gain that persists while paying its reflection from the same ledger lies strictly between them, within $c\cdot10^{-8}/23$ of $128/989$ for charge rate $c$ (\texttt{PersistentGain}). Neither comparison contains $v$, and therefore neither independently fixes the electroweak scale.

At the core, defining $M_P=m_e/\sqrt{\alpha_G}$ makes $Z_G(M_P)=1$ algebraically. Lean verifies this normalization and its uniqueness for fixed positive inputs. This remains a consistency identity rather than an independent derivation of the physical gravitational coupling.

\subsection{Joint Gauge Response}
The displayed neutral response has a massless mixed mode and a massive mode. Using its canonical normalization with the proposed VEV and the displayed uncorrected couplings gives
\begin{equation}
 78.96<m_W^{\rm response}<78.97\ {\rm GeV},\qquad
 89.57<m_Z^{\rm response}<89.58\ {\rm GeV}.
\end{equation}
These intervals use the existing electron mass-unit input and are certified by Lean. No value of $\eta$ in $[0,1]$ can close them. One interface factor $1.0362<k<1.0365$ on $\alpha$ closes both while preserving $m_W^2=(1-s)m_Z^2$, but the same factor would push the strong readout above $0.1221$ (\texttt{ScaleConsistency}). With the measured Z mass the ratio gives $80.3845<m_W<80.3846$~GeV and an implied $\Delta r\approx0.0351$, compared with about $0.0368$ from standard radiative corrections on the same inputs. This is a falsifiable difference; deriving the interface response from the substrate is the corresponding obligation.

\subsection{Status of the Constant Derivation}
The present results establish conditional resource-selection theorems, exact algebraic identities, certified numerical intervals and explicit closure residuals. They do not yet establish a unique globally closed vacuum dynamics. The remaining work is to derive the operational allocations, field response, scale dependence and spectra from the substrate while preserving the general persistence constraints. Numerical agreement of separately evaluated formulas does not replace these joint tests.
